% !TeX program = pdflatex
\documentclass[aspectratio=169,xcolor={dvipsnames,table}]{beamer}

\input{../../_en_preambulo_beamer}
% Comandos y macros estadísticas compartidas para las presentaciones Beamer en inglés (Metropolis)

% Conjuntos numéricos básicos
\newcommand{\R}{\mathbb{R}}
\newcommand{\N}{\mathbb{N}}
\newcommand{\Q}{\mathbb{Q}}
\newcommand{\Z}{\mathbb{Z}}
\newcommand{\set}[1]{\left\{ #1 \right\}}
\newcommand{\abs}[1]{\left|#1\right|}
\newcommand{\norm}[1]{\left\| #1 \right\|}

% Comandos estadísticos y probabilísticos del curso
\newcommand{\Var}{\operatorname{Var}}
\newcommand{\cov}{\operatorname{Cov}}
\newcommand{\corr}{\rho}
\newcommand{\comb}[2]{\begin{pmatrix} #1 \\ #2 \end{pmatrix}}
\newcommand{\s}{\sigma}
\newcommand{\particion}{\mathfrak{P}}
\newcommand{\card}[1]{\#\left( #1 \right)}
\newcommand{\seq}[3]{\set{#1}_{#2}^{#3}}
\newcommand{\E}{\mathbb{E}}
\newcommand{\Prob}{\mathbb{P}}

% Entornos pedagógicos y cajas en inglés académico natural (soportando tanto nombres en inglés como en español)
\newenvironment{discussion}{%
  \begin{alertblock}{Discussion Question}%
}{%
  \end{alertblock}%
}
\newenvironment{discusion}{%
  \begin{alertblock}{Discussion Question}%
}{%
  \end{alertblock}%
}

\newenvironment{reflection}{%
  \begin{alertblock}{Classroom Reflection}%
}{%
  \end{alertblock}%
}
\newenvironment{reflexion}{%
  \begin{alertblock}{Classroom Reflection}%
}{%
  \end{alertblock}%
}

\newenvironment{paradox}{%
  \begin{alertblock}{Exploratory Paradox}%
}{%
  \end{alertblock}%
}
\newenvironment{paradoja}{%
  \begin{alertblock}{Exploratory Paradox}%
}{%
  \end{alertblock}%
}

\newenvironment{motivation}{%
  \begin{exampleblock}{Motivation: Data Science in Practice}%
}{%
  \end{exampleblock}%
}
\newenvironment{motivacion}{%
  \begin{exampleblock}{Motivation: Data Science in Practice}%
}{%
  \end{exampleblock}%
}

\newenvironment{quickchallenge}{%
  \begin{block}{Quick Team Challenge}%
}{%
  \end{block}%
}
\newenvironment{retoclase}{%
  \begin{block}{Quick Team Challenge}%
}{%
  \end{block}%
}


\title{Student's t-Distribution and Small Samples}
\subtitle{Section 05.04 --- Inference on Mu with Unknown Sigma}
\author[J. Castillo Colmenares]{Juliho Castillo Colmenares}
\institute[Tec de Monterrey]{Tecnologico de Monterrey}
\date{\vspace{-1.5cm}}

\begin{document}

% SLIDE 01: Title Page
\begin{frame}[plain]
	\titlepage
\end{frame}

% SLIDE 02: Roadmap
\begin{frame}{Roadmap --- Chapter 05: Sampling Distributions}
	\vspace{-0.15cm}
	\begin{block}{Curricular Structure of Unit 4}
		\footnotesize
		\begin{enumerate}\itemsep=0.03cm
			\item \textcolor{gray}{Section 05.01: Simple Random Sampling, Mean, and Unbiased Sample Variance}
			\item \textcolor{gray}{Section 05.02: Asymptotic Central Limit Theorem}
			\item \textcolor{gray}{Section 05.03: Chi-Squared Distribution and Sample Variance}
			\item \textbf{\textcolor{TecRojo}{Section 05.04: Student's $t$-Distribution and Small Samples}} $\leftarrow$ \emph{Current focus}
			\item \textcolor{gray}{Section 05.05: Fisher-Snedecor $F$-Distribution}
		\end{enumerate}
	\end{block}
	\vspace{-0.2cm}
	\begin{alertblock}{Learning Objective}
		\scriptsize
		Formalize Student's $t$-distribution as a direct consequence of Fisher's theorem, construct exact confidence intervals for $\mu$ when $\sigma$ is unknown, and understand why small samples require an adjustment to the Normal approximation.
	\end{alertblock}
\end{frame}

% SLIDE 03: Motivation
\begin{frame}{The Price of Not Knowing $\sigma$}
	\footnotesize
	\begin{motivacion}
		In Section 05.03 we saw that $T=(\bar X-\mu)/(S/\sqrt n) = Z/\sqrt{\chi^2_{n-1}/(n-1)}$: replacing the unknown $\sigma$ with the random $S$ makes the statistic no longer exactly Normal. Today we formalize its exact distribution and practical consequences.
	\end{motivacion}
	\vspace{0.15cm}
	\pause
	\begin{itemize}\itemsep=0.1cm
		\item \textbf{Heavier tails:} $S$ introduces additional variability compared to using a fixed $\sigma$. \pause
		\item \textbf{Effect for small $n$:} the adjustment is substantial for small samples, negligible for large ones. \pause
		\item \textbf{Convergence:} as $n\to\infty$, $t_{n-1}\to N(0,1)$.
	\end{itemize}
\end{frame}

% SLIDE 03B: Definition and Properties
\begin{frame}{Definition and Properties of the $t$-Distribution}
	\vspace{-0.15cm}
	\begin{columns}[T]
		\begin{column}{0.48\textwidth}
			\begin{block}{Definition}
				\footnotesize
				For independent $Z\sim N(0,1)$ and $\chi^2_\nu$: $T=Z/\sqrt{\chi^2_\nu/\nu}\sim t_\nu$, with heavier tails than Normal.
			\end{block}
		\end{column}
		\begin{column}{0.48\textwidth}
			\begin{alertblock}{Properties}
				\footnotesize
				$E(T)=0$ ($\nu>1$); $\Var(T)=\nu/(\nu-2)$ ($\nu>2$), always $>1$; $t_\nu \xrightarrow{d} N(0,1)$ as $\nu\to\infty$.
			\end{alertblock}
		\end{column}
	\end{columns}
\end{frame}

% SLIDE 04: Confidence Interval
\begin{frame}{Confidence Interval for $\mu$ with Unknown $\sigma$}
	\vspace{-0.15cm}
	\begin{columns}[T]
		\begin{column}{0.48\textwidth}
			\begin{block}{Interval Formula}
				\footnotesize
				$\bar{X} \pm t_{n-1,\,\alpha/2}\cdot S/\sqrt{n}$. Always wider than the (incorrect) $Z$-based interval.
			\end{block}
		\end{column}
		\begin{column}{0.48\textwidth}
			\begin{alertblock}{Why It Matters More for Small $n$}
				\footnotesize
				$t_{9,0.025}\approx 2.262$ vs. $z_{0.025}=1.96$: over $15\%$ wider. The gap shrinks to negligible as $n\to\infty$.
			\end{alertblock}
		\end{column}
	\end{columns}
\end{frame}

% SLIDE 05: One-Sample t-Test
\begin{frame}{The One-Sample $t$-Test}
	\vspace{-0.15cm}
	\begin{block}{Test Statistic}
		$$t = \dfrac{\bar{X}-\mu_0}{S/\sqrt{n}} \sim t_{n-1} \text{ under } H_0: \mu=\mu_0$$
	\end{block}
	\pause
	\vspace{0.1cm}
	\begin{alertblock}{Decision Rule}
		Reject $H_0$ at level $\alpha$ if $|t| > t_{n-1,\alpha/2}$ (two-sided test). This is the classical $t$-test, standard when $n<30$ and $\sigma$ is unknown.
	\end{alertblock}
\end{frame}

% SLIDE 06: Applications
\begin{frame}{Applications in Small-Sample Inference}
	\vspace{-0.1cm}
	\begin{itemize}\itemsep=0.1cm
		\item \textbf{Quality control:} strength testing with small production batches. \pause
		\item \textbf{Clinical research:} trials with few patients; before/after comparisons (paired samples). \pause
		\item \textbf{Early A/B testing:} decisions with limited initial sample sizes. \pause
		\item \textbf{Preview 05.05:} the $F$-distribution compares variances via the ratio of two chi-squared variables, completing the $\chi^2$--$t$--$F$ trio.
	\end{itemize}
\end{frame}

% SLIDE 07: Python Lab Block 1
\begin{frame}[fragile]{Python Lab: $t$-Distribution Properties and Convergence}
	\vspace{-0.05cm}
	\scriptsize
	Verification of $\Var(T)$ and convergence of $t$ quantiles to Normal quantiles (\texttt{05.04\_student\_t\_distribution.py}):
	{\lstset{basicstyle=\fontsize{4pt}{4.8pt}\selectfont\ttfamily,aboveskip=0.1em,belowskip=0.1em}
	\lstinputlisting[firstline=17, lastline=32]{../../code/05_distribuciones_muestreo/05.04_student_t_distribution.py}}
\end{frame}

% SLIDE 08: Python Lab Block 2
\begin{frame}[fragile]{Python Lab: $t$-Based vs. $Z$-Based Intervals}
	\vspace{-0.05cm}
	\scriptsize
	Comparison of $t$-based and (incorrect) $Z$-based confidence intervals for the same dataset:
	{\lstset{basicstyle=\fontsize{4pt}{4.8pt}\selectfont\ttfamily,aboveskip=0.1em,belowskip=0.1em}
	\lstinputlisting[firstline=35, lastline=59]{../../code/05_distribuciones_muestreo/05.04_student_t_distribution.py}}
\end{frame}

% SLIDE 09: Python Lab Block 3
\begin{frame}[fragile]{Python Lab: One-Sample $t$-Test and CI Coverage}
	\vspace{-0.05cm}
	\scriptsize
	One-sample $t$-test and Monte Carlo verification of confidence interval coverage:
	{\lstset{basicstyle=\fontsize{4pt}{4.8pt}\selectfont\ttfamily,aboveskip=0.1em,belowskip=0.1em}
	\lstinputlisting[firstline=62, lastline=86]{../../code/05_distribuciones_muestreo/05.04_student_t_distribution.py}}
\end{frame}

% SLIDE 10: Python Lab Terminal Output
\begin{frame}[fragile]{Python Lab: Terminal Output}
	\vspace{-0.1cm}
	\begin{block}{}
		\fontsize{4pt}{5pt}\selectfont
		\begin{verbatim}
=== Block 1: t-Distribution Properties and Convergence ===
t_10: Var(T)=1.2500 | t_30: Var(T)=1.0714
nu=5: t=2.5706 (+31.15%) | nu=30: t=2.0423 (+4.20%) | nu=120: t=1.9799 (+1.02%)

=== Block 2: t-Based vs Z-Based Confidence Intervals ===
n=9: 95% CI (t): [46.04, 50.96] vs (z, incorrect): [46.41, 50.59] (17.7% wider)
Dataset {23,25,21,27,24}: xbar=24.00, S=2.24, 95% CI: [21.22, 26.78]

=== Block 3: One-Sample t-Test and CI Coverage ===
t=1.75, critical t_24,0.025=2.06 -> Reject H0? No
Monte Carlo CI coverage (n=10): 0.9509 (nominal: 0.9500)
		\end{verbatim}
	\end{block}
\end{frame}

% SLIDE 11: Exercise Level 1 (Statement)
\begin{frame}{In-Class Exercise --- Level 1: Fundamental (1/2)}
	\vspace{-0.1cm}
	\begin{block}{Problem 5.4.1 --- Variance of the $t$ (Statement)}
		Let $T \sim t_{10}$. Compute $\Var(T)$.
	\end{block}
	\vspace{0.15cm}
	\pause
	\begin{alertblock}{Model Setup}
		\begin{itemize}\itemsep=0.04cm
			\item Use directly $\Var(T)=\nu/(\nu-2)$.
		\end{itemize}
	\end{alertblock}
\end{frame}

% SLIDE 12: Exercise Level 1 (Solution)
\begin{frame}{In-Class Exercise --- Level 1: Fundamental (2/2)}
	\vspace{-0.05cm}
	\scriptsize
	Substitute $\nu=10$ directly: \pause
	\begin{block}{Calculation}
		$$\Var(T) = \frac{10}{10-2} = \frac{10}{8} = 1.25.$$
	\end{block}
	\vspace{0.05cm}
	\pause
	\begin{alertblock}{Interpretation}
		\scriptsize
		Since $1.25 > 1 = \Var(Z)$, $t$ with $\nu=10$ is more dispersed than the standard Normal, reflecting the additional uncertainty of estimating $\sigma$ with $S$. As $\nu$ grows, $\Var(T)\to 1$.
	\end{alertblock}
\end{frame}

% SLIDE 13: Exercise Level 2 (Statement)
\begin{frame}{In-Class Exercise --- Level 2: Operational (1/2)}
	\vspace{-0.1cm}
	\begin{block}{Problem 5.4.5 --- One-Sample $t$-Test (Statement)}
		It is claimed that $\mu_0=100$. A sample of $n=25$ yields $\bar X=104.2$ and $S=12$. Compute $t=(\bar X-\mu_0)/(S/\sqrt n)$ and decide whether to reject $H_0:\mu=100$ at $\alpha=0.05$, using $t_{24,0.025}\approx 2.064$.
	\end{block}
	\vspace{0.15cm}
	\pause
	\begin{alertblock}{Strategy}
		\scriptsize
		Compare $|t|$ against the critical value $t_{24,0.025}$.
	\end{alertblock}
\end{frame}

% SLIDE 14: Exercise Level 2 (Solution)
\begin{frame}{In-Class Exercise --- Level 2: Operational (2/2)}
	\vspace{-0.05cm}
	\scriptsize
	Compute the statistic: \pause
	\begin{block}{Calculation}
		$$t = \frac{104.2-100}{12/5} = \frac{4.2}{2.4} = 1.75$$
	\end{block}
	\vspace{0.05cm}
	\pause
	\begin{alertblock}{Decision}
		\scriptsize
		Since $|1.75| < 2.064$, we \textbf{fail to reject} $H_0:\mu=100$ at $\alpha=0.05$: the sample does not provide sufficient evidence that the mean differs from $100$.
	\end{alertblock}
\end{frame}

% SLIDE 15: Exercise Level 3 (Statement)
\begin{frame}{In-Class Exercise --- Level 3: Analytical (1/2)}
	\vspace{-0.1cm}
	\begin{block}{Problem 5.4.7 --- Deriving $\Var(T)$ (Statement)}
		Given that $E(1/\chi^2_\nu)=1/(\nu-2)$ for $\nu>2$, that $Z$ and $\chi^2_\nu$ are independent, and that $E(Z^2)=1$, derive $\Var(T)=\nu/(\nu-2)$ for $T=Z/\sqrt{\chi^2_\nu/\nu}$.
	\end{block}
	\vspace{0.1cm}
	\pause
	\begin{alertblock}{Strategy}
		\scriptsize
		Since $E(T)=0$, $\Var(T)=E(T^2)$. Write $T^2$ in terms of $Z^2$ and $\chi^2_\nu$, then use independence to factor the expectation.
	\end{alertblock}
\end{frame}

% SLIDE 16: Exercise Level 3 (Solution)
\begin{frame}{In-Class Exercise --- Level 3: Analytical (2/2)}
	\vspace{-0.05cm}
	\scriptsize
	Factor using independence: \pause
	\begin{block}{Derivation}
		\scriptsize
		\begin{align*}
			T^2 = \frac{Z^2}{\chi^2_\nu/\nu} = \nu\frac{Z^2}{\chi^2_\nu} \implies
			E(T^2) = \nu\, E(Z^2)\, E\!\left(\frac{1}{\chi^2_\nu}\right) = \nu\cdot 1\cdot\frac{1}{\nu-2} = \frac{\nu}{\nu-2}. \quad \qedhere
		\end{align*}
	\end{block}
	\vspace{0.05cm}
	\pause
	\begin{alertblock}{Interpretation}
		\scriptsize
		Since $E(T)=0$, $\Var(T)=E(T^2)=\nu/(\nu-2)$, confirming the formula. The independence of $Z$ and $\chi^2_\nu$ is essential to factor the expectation of the product.
	\end{alertblock}
\end{frame}

% SLIDE 17: Exercise Level 4 (Statement)
\begin{frame}{In-Class Exercise --- Level 4: Challenging (1/2)}
	\vspace{-0.1cm}
	\begin{block}{Problem 5.4.9 --- Paired Samples (Statement)}
		In a diet study, $n=8$ subjects are weighed before and after. The mean difference is $\bar d=-3.2$ kg with $S_d=2.1$ kg. Construct a $95\%$ CI for the mean difference using $t_{7,0.025}\approx 2.365$, and determine whether there is evidence the diet reduces weight.
	\end{block}
	\vspace{0.15cm}
	\pause
	\begin{alertblock}{Strategy}
		\scriptsize
		Treat the differences $d_i$ as a single sample; if the interval excludes zero, there is evidence of a real effect.
	\end{alertblock}
\end{frame}

% SLIDE 18: Exercise Level 4 (Solution)
\begin{frame}{In-Class Exercise --- Level 4: Challenging (2/2)}
	\vspace{-0.05cm}
	\scriptsize
	Compute the standard error and margin: \pause
	\begin{block}{Calculation}
		\scriptsize
		\begin{align*}
			\text{SE} = \frac{2.1}{\sqrt 8} \approx 0.7425, \qquad \text{margin} = 2.365\cdot 0.7425 \approx 1.756.
		\end{align*}
		\begin{align*}
			\text{95\% CI} = -3.2 \pm 1.756 = [-4.956,\ -1.444].
		\end{align*}
	\end{block}
	\vspace{0.05cm}
	\pause
	\begin{alertblock}{Interpretation}
		\scriptsize
		Since the entire interval is negative (excludes zero), there is statistical evidence that the diet reduces average weight by roughly $1.44$ to $4.96$ kg, with $95\%$ confidence.
	\end{alertblock}
\end{frame}

% SLIDE 19: Closing and Chapter Perspective
\begin{frame}{Closing Section and Chapter Perspective}
	\vspace{-0.2cm}
	\begin{block}{Synthesis: Student's $t$-Distribution}
		\scriptsize
		The $t$-distribution solves the most common practical problem in statistical inference: estimating $\mu$ when $\sigma$ is unknown. Its construction --- a Normal divided by the square root of an independent chi-squared --- is not arbitrary, but a direct consequence of Fisher's theorem from the previous session.
	\end{block}
	\vspace{0.05cm}
	\pause
	\begin{alertblock}{Key Lessons and Next Step}
		\begin{itemize}\itemsep=0.05cm
			\item \textbf{Definition:} $T=Z/\sqrt{\chi^2_\nu/\nu}$; heavier tails than Normal, $\Var(T)=\nu/(\nu-2)$.
			\item \textbf{CI for $\mu$:} $\bar X \pm t_{n-1,\alpha/2}\cdot S/\sqrt n$, always wider than the $Z$-based interval.
			\item \textbf{Convergence:} $t_\nu \xrightarrow{d} N(0,1)$ as $\nu\to\infty$; the effect is most noticeable for small samples.
			\item \textbf{Next Section:} \textcolor{TecRojo}{Fisher-Snedecor $F$-Distribution}.
		\end{itemize}
	\end{alertblock}
\end{frame}

\end{document}
